\chapter{Troubleshooting}\label{ch:troubleshooting}

\app\ explains a problem in a message rather than with a traceback, and names the
elements involved. This chapter lists the messages you are most likely to meet, grouped
by when they appear, with what to do about each one.

\section{Validation: before any study runs}

These are checked on every run. An \emph{error} stops the run, and a \emph{warning} lets
it go on.

\begin{reftable}{@{}>{\raggedright\arraybackslash}p{6.2cm}X@{}}
\toprule
\sffamily\bfseries Message & \sffamily\bfseries What to do \\ \midrule
\emph{No base reference bus is set.} & \menu{Network > Base settings...}: pick the bus
  where you know the base voltage, and give that voltage. \\
\emph{The base voltage at the reference bus must be positive.} & The same dialog: set
  \gui{Base voltage there}. \\
\emph{X is not connected to a bus / at both ends, so the studies ignore it.}
  (warning) & Wire the terminal (\cref{sec:connecting}), or delete the element. Every
  result repeats this warning while it lasts. \\
\emph{X connects bus 'a' to itself.} & Reconnect one end of the branch to another bus. \\
\emph{X: rated voltages must be positive. / rated power must be positive.} & Fill in the
  transformer's ratings. \\
\emph{X has zero impedance.} (warning; an error when a study converts it) & Give the
  transformer, line or impedance a non-zero $r$ or $x$. A zero impedance is a short
  circuit between two buses: merge them instead. \\
\emph{X: cos phi must be in (0, 1].} & Correct the load's power factor, or give it as
  $P$ and $Q$. \\
\emph{X: the voltage setpoint must be positive. / Q min is above Q max.} & Correct the
  generator's or SVC's setpoint or limits. \\
\emph{The network has N separate islands.} (warning) & Usually intended, for example a
  case that disconnects part of the network. The network studies keep the islands that
  contain a source. \\
\bottomrule
\end{reftable}

\section{Bases}

\begin{reftable}{@{}>{\raggedright\arraybackslash}p{6.2cm}X@{}}
\toprule
\sffamily\bfseries Message & \sffamily\bfseries What to do \\ \midrule
\emph{Bus 'x' already has U\_base = 12 kV but is also reached with 12.6 kV through 'T3'.
  Check the transformer ratios around this loop.} & The rated ratios around a loop do not
  multiply to one. Correct a ratio, or pin the base you want on one bus
  (\cref{ch:bases}). \\
\emph{Bus 'x' cannot be reached from the base reference bus, so it has no per-unit base.}
  & Connect the bus to the rest of the network, pin its base, or take it out of
  service. \\
\emph{X spans a change of voltage base (40 kV vs 12 kV). A line cannot change the base;
  use a transformer.} & A line joins two zones. Use a transformer between them, or
  correct the bases. \\
\bottomrule
\end{reftable}

\section{Studies}

\begin{reftable}{@{}>{\raggedright\arraybackslash}p{6.2cm}X@{}}
\toprule
\sffamily\bfseries Message & \sffamily\bfseries What to do \\ \midrule
\emph{There is no source in the network, so the bus voltages are not defined. Add an
  external grid.} & The linear, Thevenin, fault and unbalanced-load studies need an
  external grid. A generator only acts in the load flow. \\
\emph{X has no zero-sequence data. Set 'Z0 data' in its Zero sequence properties.} &
  Give the line its zero-sequence data, as factors or per-km values. The data manager's
  \gui{Line / Cable} tab sets it for many lines at once. \\
\emph{Bus 'x' holds an ideal source, so its impedance to ground is zero by construction.
  Give the source a finite impedance.} & Give the external grid an impedance (\gui{Z
  Thevenin}, or a fault level), or study another bus. \\
\emph{Bus 'x' is not energized in this case.} & The case takes out of service the path
  from the sources to that bus. Check its \gui{Network variant}. \\
\emph{The condition at bus 'x' cannot be solved: it short-circuits an ideal source or
  leaves the currents undetermined.} & A fault at, or next to, an ideal source. Give the
  source an impedance. \\
\emph{The two ports must be different buses.} (S1a) & Choose two different port buses. \\
\emph{Nothing is connected to bus 'x', so its equivalent impedance is infinite.} (S1b) &
  Choose a bus with something connected, or connect it. \\
\emph{The network cannot be solved: some buses are floating (no path to ground and no
  source).} & A part of the network has no shunt and no source: add the missing
  connection, or take the floating part out of service. \\
\bottomrule
\end{reftable}

\section{Load flow}

\begin{reftable}{@{}>{\raggedright\arraybackslash}p{6.2cm}X@{}}
\toprule
\sffamily\bfseries Message & \sffamily\bfseries What to do \\ \midrule
\emph{The island with buses \dots\ has no slack bus: connect an external grid or set a
  generator's control to slack.} & Every island needs one slack bus. \\
\emph{The island with buses \dots\ has 2 slack buses; a load flow needs exactly one per
  island.} & Set all but one to \code{pv}, or clear \gui{Slack} on extra external
  grids. \\
\emph{Conflicting voltage setpoints at bus x: G7 and G7b ask for different values.} &
  Give every voltage-controlling device on a bus the same \gui{U set}. \\
\emph{The load flow did not converge.} (warning) & Read the iteration log: a mismatch
  that grows points at the bus named in the \gui{at bus} column. Typical causes are a load
  beyond what the network can carry, a voltage setpoint far from 1~pu, or wrong units,
  such as kW typed as MW. Raise \gui{Maximum iterations} only when the mismatch was still
  falling. \\
\emph{Bus x reached its max reactive limit and becomes a PQ bus.} (warning, with
  \gui{Enforce reactive limits}) & Information: the bus can no longer hold its voltage.
  Its voltage is now a result. \\
\bottomrule
\end{reftable}

\section{Study cases}

\begin{reftable}{@{}>{\raggedright\arraybackslash}p{6.2cm}X@{}}
\toprule
\sffamily\bfseries Message & \sffamily\bfseries What to do \\ \midrule
\emph{Override 'x': write it as COMPONENT\_ID.field, e.g. LD4.p\_mw.} & The dialog refuses
  an override whose element or field does not exist. Use the element's id, not its name,
  and a field name from \cref{ch:components}. \\
\emph{Case 'c': override 'x' names something that does not exist, so it is ignored.}
  (warning) & An element the case refers to was deleted or renamed. Edit the case. \\
\emph{Unknown name 'B'. Define it as a variable or load it from your script.} & An
  expression uses a variable that does not exist. Add it in the \gui{Variables}
  panel. \\
\bottomrule
\end{reftable}

\section{Installation and display}

\begin{description}
  \item[Qt fails to load] with \emph{``the specified procedure could not be found''}. On
    Windows with Anaconda, Anaconda's ICU library shadows the one Qt expects. Start
    \app\ with \file{run\_gui.py}, which imports \code{psa.gui} first; in your own
    scripts, import \code{psa.gui} before PySide6.
  \item[The window layout looks wrong] Panels can be dragged back into place, or shown
    again from \gui{View}. To return to the default layout, close \app\ and delete the
    \code{windowState} value under \file{HKCU\textbackslash Software\textbackslash
    PowerLab\textbackslash PowerLab} in the registry.
  \item[The results do not change after an edit] Results are not recomputed by
    themselves. When the tabs carry an asterisk, press \keys{F5}.
  \item[A pasted element is unconnected] Its bus was not copied with it
    (\cref{sec:clipboard}). Wire it, or copy the bus too.
\end{description}

\section{Limitations of version 1.0}\label{sec:limitations}

\begin{itemize}
  \item One diagram page per project. No breakers or switches as components (take
    elements out of service instead). No three-winding transformers.
  \item Transformer taps are stored but not modelled (an off-nominal tap only gives a
    warning), and there are no phase shifters. Phase shifts of vector groups are not
    modelled in the sequence networks.
  \item Short circuits are computed from the linear pre-fault state, as in the course,
    without the voltage factors of IEC~60909.
  \item The load flow has one slack per island (no distributed slack), and a PU bus that
    reaches a reactive limit stays PQ (no switching back).
  \item The data manager edits one cell at a time.
\end{itemize}
